\documentclass[11pt]{article}

\usepackage[margin=0.9in]{geometry}
\usepackage{amsmath,amssymb,amsthm,mathtools}
\usepackage{graphicx}
\usepackage{xcolor}
\usepackage{microtype}
\usepackage[hidelinks]{hyperref}

\newtheorem{theorem}{Theorem}
\newtheorem{corollary}{Corollary}
\newcommand{\R}{\mathbb R}
\newcommand{\Kclass}{\mathcal K^n_{+,0}}
\newcommand{\vol}{\operatorname{Vol}}

\title{$p$-Affine Surface Areas Generate Centro-Affine Curvature Integrals\\
\large A full positive answer to the integral-generation conjecture in arXiv:1208.0783}
\author{Automated mathematics run \texttt{fa\_banach\_001}, lane 14}
\date{22 July 2026}

\begin{document}
\maketitle

\begin{center}
\fcolorbox{black}{gray!8}{\parbox{0.88\textwidth}{
\textbf{Status.} Candidate full solution, likely valid, pending human review.
The theorem answers the precise conjecture concerning integrals of continuous
functions of positive centro-affine curvature.
}}
\end{center}

\section{Source conjecture}

Let $\Kclass$ denote the class of $C^2_+$ convex bodies in $\R^n$ containing
the origin in their interior. For $K\in\Kclass$, write $\kappa_0$ for its
centro-affine curvature and $d\mu_K$ for cone measure on $\partial K$. Then
$\kappa_0$ is continuous and strictly positive, and
\begin{equation}
 \mu_K(\partial K)=n\vol(K).
 \label{eq:mass}
\end{equation}
The source \cite{Stancu} uses, for $p\in\R\setminus\{-n\}$,
\begin{equation}
 \Omega_p(K)=\int_{\partial K}
 \kappa_0^{\frac{p}{n+p}}\,d\mu_K.
 \label{eq:omega}
\end{equation}
On p.~12 it conjectures that the $p$-affine surface areas, algebraic
operations, and closure generate every invariant of the form
\begin{equation}
 I_\phi(K)=\int_{\partial K}\phi(\kappa_0)\,d\mu_K.
 \label{eq:Iphi}
\end{equation}
The complete source paragraph is reproduced here.

\begin{figure}[htbp]
 \centering
 \includegraphics[width=\textwidth]{figures/open_problem_crop.png}
 \caption{The integral-generation conjecture on p.~12 of the source PDF.}
 \label{fig:source}
\end{figure}

Because the source does not specify a topology for the closure, we prove both
pointwise closure on all of $\Kclass$ and uniform closure on every natural
curvature/volume-controlled stratum.

\section{Proof intuition}

Formula \eqref{eq:omega} says that $p$-affine surface areas are moments of
$\kappa_0$. The parameter change $t=p/(n+p)$ attains every real exponent
except $1$. In particular, choosing
\[
 p_k=-\frac{2kn}{2k-1}
\]
makes $t=2k$. Hence a finite linear combination of the corresponding
$\Omega_{p_k}$ is exactly the integral of a polynomial in $\kappa_0^2$.

Every continuous $\phi$ on $(0,\infty)$ can be written as
$\phi(x)=\psi(x^2)$ with $\psi(y)=\phi(\sqrt y)$ continuous. Approximate
$\psi$ locally uniformly by ordinary polynomials. Each fixed positive
continuous curvature function has compact range inside $(0,\infty)$, so the
polynomial approximation can be integrated. This already uses only the
linear span of the $p$-affine surface areas; products of invariants are not
needed.

\section{Main theorem}

For $a,b,V>0$ with $a\leq b$, define
\[
 \Kclass(a,b,V)=\left\{K\in\Kclass:
 a\leq\kappa_0\leq b\text{ on }\partial K,
 \ \vol(K)\leq V\right\}.
\]

\begin{theorem}[Generation by a fixed countable parameter set]
Let $n\geq2$ and $\phi\in C((0,\infty))$. Put
\[
 p_0=0,
 \qquad
 p_k=-\frac{2kn}{2k-1}\quad(k\geq1).
\]
There are finite linear combinations
\[
 F_m=\sum_{k=0}^{d_m}c_{m,k}\Omega_{p_k},
\]
whose coefficients depend only on $\phi$ and $m$, such that
\begin{equation}
 F_m(K)\longrightarrow I_\phi(K)
 \qquad\text{for every }K\in\Kclass.
 \label{eq:pointwise}
\end{equation}
Moreover, for every $0<a\leq b<\infty$ and $V<\infty$, the convergence is
uniform on $\Kclass(a,b,V)$.
\end{theorem}

\begin{proof}
Set $\psi(y)=\phi(\sqrt y)$ for $y>0$. For every integer $m\geq1$, the
classical polynomial approximation theorem \cite[Chapter 7]{Rudin} gives a polynomial $Q_m$ such
that
\begin{equation}
 \sup_{y\in[m^{-2},m^2]}|Q_m(y)-\psi(y)|<\frac1m.
 \label{eq:approx}
\end{equation}
Thus $Q_m(x^2)\to\phi(x)$ locally uniformly for $x\in(0,\infty)$.

Write $Q_m(y)=\sum_{k=0}^{d_m}c_{m,k}y^k$. For $k\geq1$, direct algebra gives
\begin{equation}
 p_k<-n,
 \qquad
 \frac{p_k}{n+p_k}=2k.
 \label{eq:parameter}
\end{equation}
Thus $p_k$ avoids the sole singular parameter $-n$, and
\eqref{eq:omega} gives
\[
 \Omega_{p_k}(K)=\int_{\partial K}\kappa_0^{2k}\,d\mu_K.
\]
For $k=0$, $p_0=0$ and
$\Omega_0(K)=\int_{\partial K}1\,d\mu_K$. Consequently,
\begin{align}
 F_m(K)
 &:=\sum_{k=0}^{d_m}c_{m,k}\Omega_{p_k}(K)\\
 &=\int_{\partial K}Q_m(\kappa_0^2)\,d\mu_K.
 \label{eq:realization}
\end{align}

Fix $K\in\Kclass$. Compactness of $\partial K$ and strict positivity of
$\kappa_0$ give constants $0<a_K\leq b_K<\infty$ with
$a_K\leq\kappa_0\leq b_K$. Local uniform convergence therefore implies
\[
 |F_m(K)-I_\phi(K)|
 \leq n\vol(K)\sup_{x\in[a_K,b_K]}
 |Q_m(x^2)-\phi(x)|\longrightarrow0,
\]
where we used \eqref{eq:mass}. This proves \eqref{eq:pointwise}.

If $K\in\Kclass(a,b,V)$, the same estimate is bounded by
\[
 nV\sup_{x\in[a,b]}|Q_m(x^2)-\phi(x)|,
\]
which tends to zero independently of $K$. This proves the uniform assertion.
\end{proof}

\begin{corollary}[Normalized form]
Let
\[
 \widehat\Omega_p(K)=\frac{\Omega_p(K)}{\Omega_0(K)},
 \qquad
 \widehat I_\phi(K)=\frac{I_\phi(K)}{\Omega_0(K)}.
\]
Then
\[
 \sum_{k=0}^{d_m}c_{m,k}\widehat\Omega_{p_k}(K)
 \longrightarrow\widehat I_\phi(K)
\]
pointwise on $\Kclass$, and uniformly over all $K$ satisfying only
$a\leq\kappa_0\leq b$, with no volume bound.
\end{corollary}

\begin{proof}
The normalized cone measure
$d\nu_K=d\mu_K/\Omega_0(K)$ is a probability measure. Divide
\eqref{eq:realization} by $\Omega_0(K)$; the approximation error is then at
most $\sup_{[a,b]}|Q_m(x^2)-\phi(x)|$.
\end{proof}

\section{Consequences, verification, and scope}

The theorem places every $I_\phi$ in the pointwise closure of the linear span
of $\{\Omega_{p_k}:k\geq0\}$ and hence in the closure of the algebra generated
by all $p$-affine surface areas. The parameters form a fixed countable set and
satisfy $p_k\uparrow-n$ from below. The normalized corollary directly matches
the normalized wording in the source abstract.

The exact-arithmetic script \texttt{verify\_parameter\_map.py} checks
\eqref{eq:parameter} for 570 pairs. It is a bookkeeping check, not a substitute
for the proof. The negative parameters are harmless because $\kappa_0$ is
positive and bounded away from zero.

The four run indexes and two bounded web-search rounds used the exact
conjecture wording and close mathematical variants. No matching resolution
was found; novelty confidence is moderate because the search was not exhaustive.

The result answers the precise p.~12 conjecture \eqref{eq:Iphi}, not the
abstract's broader informal phrasing about arbitrary $SL(n)$-invariants.
Review should focus on the intended closure topology and the admitted range
$p\in\R\setminus\{-n\}$.

\begin{thebibliography}{9}
\bibitem{Stancu}
A. Stancu,
\emph{Some Affine Invariants Revisited},
in \emph{Asymptotic Geometric Analysis}, Fields Institute Communications 68
(2013), 343--359;
arXiv:\href{https://arxiv.org/abs/1208.0783}{1208.0783}.

\bibitem{Rudin}
W. Rudin,
\emph{Principles of Mathematical Analysis}, 3rd ed.,
McGraw--Hill, 1976, Chapter 7.
\end{thebibliography}

\end{document}
